\chapter{Historia del desarrollo (I): del estudio único a un protocolo con integridad temporal}
\label{chap:desarrollo-metodo}

Los capítulos anteriores presentan el protocolo como si hubiera existido desde el principio. No fue
así. El diseño que se ha descrito es el resultado de una secuencia de correcciones, y varias de
ellas nacieron de descubrir que una medición previa era inválida. Este capítulo, dividido en dos
partes, documenta esa trayectoria a partir de la bitácora del proyecto. Esta primera mitad cubre el
paso de una arquitectura con demasiadas rutas a un único Model Study, los defectos que afectaron a
la metodología de contratos y a la selección secuencial, y la ponderación por recencia del
meta-agente. La segunda mitad, tras el Capítulo~\ref{chap:diseno-experimental}, cubre el rediseño de
la cartera, el colapso de la calibración isotónica y la advertencia de reproducibilidad.

Hay dos razones para incluir esta historia en lugar de presentar sólo el resultado final. La primera
es de honestidad: algunos de los defectos que se describen afectaron a cifras que llegaron a
documentarse como resultados, y omitirlos daría una imagen de linealidad que no corresponde a los
hechos. La segunda es que los propios defectos son evidencia. Un sistema que detecta que una de sus
puertas estadísticas penalizaba a los candidatos superiores, y que lo detecta porque ningún retador
resultaba elegible pese a dominar, está mostrando algo sobre su capacidad de rechazar estados
inválidos. Las lecciones son además transferibles: casi ninguna es específica de este dominio.

\section{El estudio como unidad científica}

El diseño original separaba la investigación en dos fases, una exploratoria y otra confirmatoria,
con hipótesis formuladas entre ambas. Se eliminó en julio de 2026 y se sustituyó por un único
\textit{Model Study} automático.

El motivo inmediato fue operativo. Un estudio iniciado el 24 de julio completó un ajuste costoso y
falló después, al consultar un campo de configuración que ya había sido eliminado del código. El
error vivía sólo en memoria, y el proceso desapareció sin dejar reconciliación posible: no había
forma de saber, desde los artefactos, qué se había llegado a calcular. El problema de fondo era que
la arquitectura admitía demasiadas rutas --- escenarios, experimentos, ejecuciones sueltas ---, y
cada ruta era una oportunidad para que un estado parcial quedara sin registrar.

La reconstrucción impuso una regla que el resto del trabajo hereda: existe un único tipo de unidad
científica, se ejecuta entera o falla, y todo lo que produce queda persistido con identidad propia.
Los estados posibles de una ejecución se hicieron explícitos, se añadió un latido para detectar
procesos muertos y una reconciliación por identificador de proceso al arrancar. La consecuencia
práctica para este documento es que cada cifra del Capítulo~\ref{chap:resultados} puede rastrearse
hasta un fichero concreto de un estudio concreto.

Durante la validación de esa reconstrucción aparecieron ocho incidencias que conviene mencionar
porque ilustran cuántas formas hay de que una ejecución «funcione» sin ser válida. La más
instructiva es la segunda: el primer \textit{smoke} técnicamente exitoso produjo puntuaciones
constantes. La causa era que el mínimo de cincuenta observaciones por hoja impedía que los árboles
dividieran un conjunto de sesenta y cinco filas, de modo que el modelo devolvía siempre el mismo
valor. El proceso terminó sin error, escribió sus artefactos y se declaró correcto. No se aceptó
como validación. Es un recordatorio de que la ausencia de excepción no es evidencia de corrección.

\section{Variables declaradas que no se calculaban}

El hallazgo científicamente más relevante de la limpieza de julio fue de otra naturaleza. El
catálogo de \textit{features} declaraba dos bloques de \textit{momentum} multi-horizonte, presentes
en todos los conjuntos de variables reales. Pero esas columnas sólo se calculaban si dos
interruptores concretos estaban activados, y esos dos interruptores \emph{nunca existieron en el
catálogo cerrado}: no eran alcanzables desde ningún estudio real.

La consecuencia es que un estudio previamente documentado se ejecutó sin las columnas de
\textit{momentum} multi-horizonte que su propia configuración declaraba usar. Nada falló, ninguna
prueba se puso en rojo y los resultados eran perfectamente plausibles. El defecto era invisible
porque el sistema declaraba una cosa y computaba otra, y ambas descripciones eran internamente
coherentes.

La lección --- declarar no es computar --- motivó una comprobación que hoy forma parte de la
batería de contratos: que el catálogo declarado y el conjunto de columnas efectivamente calculadas
coincidan. En la misma limpieza se retiró un sesgo de régimen fijado por constantes en el código, no
aprendido de los datos, que contradecía frontalmente el propósito del trabajo.

\section{Cuando la puerta estadística penalizaba al mejor}

La corrección de validez de finales de julio es la más extensa del registro: doce defectos
identificados y medidos, recogidos en la Tabla~\ref{tab:defectos}.

\begin{table}[H]
  \centering
  \footnotesize
  \caption{Defectos de validez detectados y efecto medido de su corrección.}
  \label{tab:defectos}
  % Fuente: docs/bitacora.md, entrada 2026-07-28 «Corrección de validez, alfa neto y limpieza».
% Tabla transcrita del registro de desarrollo, no generada por script: su origen es la
% bitácora del proyecto y no un artefacto del Study.
\begin{tabularx}{\textwidth}{XX}
\toprule
Defecto detectado & Efecto medido de la corrección \\
\midrule
La puerta de no inferioridad penalizaba a los candidatos \textbf{superiores}: cuanto mayor la
diferencia, más ancho el intervalo &
\texttt{feature\_preset} pasa de \texttt{core} (Rank-IC 0,0730) a \texttt{all} (0,0958), con
ventaja pareada $+0{,}0216$. Ningún retador era elegible pese a dominar \\
\addlinespace
\texttt{market\_regime\_feature} se decidió sobre una ventaja de 0,00112, por debajo del ruido &
Pasa de \texttt{True} a \texttt{False} por simplicidad \\
\addlinespace
\texttt{meta\_method} se decidió sobre una ventaja de 0,00033 &
Se mantiene, pero registrado como empate técnico y no como victoria \\
\addlinespace
Emparejamiento por cadena de fecha con rejillas desplazadas por \texttt{execution\_lag\_days} &
Se empareja por periodo mensual; sin bloque completo común el resultado se marca no aplicable en
vez de devolver \texttt{ci\_low = 0,0} \\
\addlinespace
\texttt{evaluation\_key} no incluía el hash del dataset &
La misma clave aparecía con CAGR 0,1468 y 0,1692. Ahora la clave separa datasets \\
\addlinespace
Seis factores de precio se inyectaban en el agente \textit{momentum} fuera de todo condicional &
La ablación \texttt{fundamental} deja de recibir información de precio y mide lo que declara \\
\addlinespace
Dos definiciones incompatibles de \textit{information ratio} bajo el mismo nombre &
Una sola, anualizada \\
\addlinespace
CAGR, \textit{drawdown} y alfa mezclaban la era reservada con la de selección &
Métricas segmentadas: selección, confirmación y curva completa \\
\addlinespace
\texttt{geometric\_excess\_return} era una resta de CAGR &
Cociente de acumulados \\
\addlinespace
Una posición sin precio se marcaba plana &
Convención de exclusión tipo CRSP ($-30\,\%$) y liquidación contra efectivo \\
\addlinespace
\texttt{subsample=0.8} sin \texttt{subsample\_freq}: el \textit{bagging} nunca se activaba &
\texttt{subsample\_freq=1} \\
\addlinespace
Nulo de carteras aleatorias con percentil 95 de CAGR del 107\,\% &
Exige cobertura anual completa, aplica la guarda de datos y paga los mismos costes \\
\bottomrule
\end{tabularx}

  \caption*{Fuente: \artefacto{docs/bitacora.md}, entrada de 2026-07-28.}
\end{table}

El primero de la lista merece detalle porque afecta al corazón del protocolo. La regla de selección
admite un candidato si demuestra no ser inferior al incumbente, y esa no inferioridad se evaluaba
con un intervalo de confianza sobre la diferencia pareada. El defecto era que el intervalo se
ensanchaba con la magnitud de la diferencia: cuanto \emph{mejor} era un retador, más ancho su
intervalo y más difícil superar el umbral. La puerta penalizaba exactamente a los candidatos que
debía premiar.

El síntoma que lo delató no fue un error, sino una regularidad sospechosa: ningún retador resultaba
elegible, en ninguna variable, pese a que algunos dominaban con claridad. Al corregirlo,
\texttt{feature\_preset} pasó de \texttt{core} (Rank-IC 0,0730) a \texttt{all} (0,0958), con una
ventaja pareada de $+0{,}0216$. Es decir, el sistema había estado descartando una mejora
sustancial y real por un defecto en su criterio de admisión.

Dos defectos más de esa lista tienen implicaciones directas para la lectura del
Capítulo~\ref{chap:resultados}. La clave de evaluación no incluía el identificador del conjunto de
datos, de modo que la misma clave llegó a aparecer asociada a CAGR 0,1468 y 0,1692: dos resultados
distintos indexados como si fueran el mismo. Y el exceso geométrico se calculaba como una resta de
CAGR en lugar de un cociente de acumulados, lo que introduce un sesgo que crece con el horizonte.
Ambas correcciones cambian cifras publicadas, y por eso los estudios anteriores a esta versión del
catálogo quedaron derogados como evidencia.

De la misma revisión salió una observación que motivó introducir un conjunto de semillas: el Rank-IC
variaba apenas $\pm 0{,}001$ entre inicializaciones, pero el exceso sobre el índice \emph{cambiaba
de signo} (una semilla daba $-0{,}51$ puntos y otra $+3{,}11$). La conclusión predictiva era
estable; la económica, no. Esa asimetría es la que la Figura~\ref{fig:semillas} documenta hoy con
tres semillas.

\section{Una ventana que no olvidaba el régimen antiguo}

En julio de 2026 se observó un comportamiento que no encajaba: posiciones con ochenta y muchos meses
de antigüedad y percentil actual muy bajo que nunca se vendían. El diagnóstico encadenó dos
problemas distintos; este apartado documenta el primero, de naturaleza predictiva, y el
Capítulo~\ref{chap:desarrollo-cartera} documenta el segundo, de naturaleza estadística.

El primer problema era que el meta-agente no olvidaba el régimen antiguo. Su Rank-IC acumuló siete
trimestres negativos consecutivos, una racha más larga que la de la crisis sanitaria, y los cinco
agentes se invirtieron a la vez --- un patrón compatible con una rotación de factores, no con un
error de signo. El ajuste usaba una ventana rectangular: una cohorte de cuatro años atrás pesaba lo
mismo que la más reciente. Este mismo patrón fue el que motivó introducir la ponderación por
recencia como variable del catálogo, cuya evidencia se documenta en el
Capítulo~\ref{chap:agentes}.

\section{Umbrales anuales y comparabilidad}

Un cambio aparentemente menor tiene consecuencias sobre todo lo demás. Los umbrales de la cartera se
expresaban en puntos básicos referidos al horizonte, no al año. Con esa convención, 250 puntos
básicos en un horizonte de tres meses equivalen aproximadamente a un 10\,\% anual, mientras que en
uno de doce meses equivalen a un 2,5\,\%: el mismo número significaba exigencias muy distintas según
otra variable del catálogo.

La corrección fue expresar todos los umbrales en puntos básicos anuales y convertirlos al horizonte
mediante composición geométrica. El prorrateo lineal se rechazó de forma deliberada: es el mismo
atajo aritmético que ya había sido corregido en el CAGR y en el \textit{information ratio}, y
reintroducirlo por comodidad habría sido inconsistente.

\section{El libro de versiones del catálogo}

Cada uno de estos cambios rompe la comparabilidad con los estudios anteriores, y el proyecto lleva
un registro explícito de esas rupturas mediante una versión de catálogo. La
Tabla~\ref{tab:versiones} lo resume junto con el crecimiento de la batería de contratos.

\begin{table}[H]
  \centering
  \footnotesize
  \caption{Versiones del catálogo cerrado y crecimiento de la batería de pruebas.}
  \label{tab:versiones}
  % Fuente: docs/bitacora.md (entradas 2026-07-25 a 2026-08-04) y module/studies/catalog.py.
% Transcrita del registro de desarrollo, no generada por script.
\begin{tabularx}{\textwidth}{clXr}
\toprule
Versión & Fecha & Qué cambia y por qué rompe la comparabilidad & Tests \\
\midrule
2 & 2026-07-25 &
Reconstrucción a Model Study único: se elimina el protocolo exploratorio--confirmatorio y se retiran
las variantes de meta con sesgo no aprendido (\texttt{regime}, \texttt{rank\_ic}) &
15 \\
\addlinespace
3 & 2026-07-28 &
Corrección de doce defectos de validez; los umbrales de cartera pasan de percentiles a puntos
básicos y \texttt{sizing\_mode} de \texttt{score\_linear} a \texttt{alpha\_proportional} &
--- \\
\addlinespace
4 & 2026-07-28 &
Rediseño de la gestión de cartera bajo el principio de que toda venta necesita un destino mejor;
\texttt{max\_cash\_weight} entra con rejilla $(0\,\%,\,10\,\%,\,25\,\%)$ &
45 \\
\addlinespace
5 & 2026-07-28 &
Umbrales expresados en puntos básicos anuales con conversión geométrica; entran
\texttt{minimum\_holding\_period} y \texttt{price\_only\_sell\_only} &
52 \\
\addlinespace
6 & 2026-08-04 &
\texttt{cash\_policy} se elimina y \texttt{max\_cash\_weight} gobierna sola el efectivo; se corrige
el reparto que dejaba efectivo no declarado y entra \texttt{coverage\_percentile\_floor} &
75 \\
\bottomrule
\end{tabularx}

  \caption*{Fuente: \artefacto{docs/bitacora.md} y \artefacto{module/studies/catalog.py}.}
\end{table}

El número de contratos pasó de 15 a 75 a lo largo del desarrollo, y casi todos los añadidos
proceden de un defecto concreto: cada corrección de esta lista dejó tras de sí una prueba que
impide su reaparición. Esa es, probablemente, la parte más reutilizable del trabajo.

La segunda mitad de esta historia --- el rediseño de la cartera, el colapso de la calibración
isotónica y la advertencia de reproducibilidad que condiciona toda cifra de este documento --- se
retoma en el Capítulo~\ref{chap:desarrollo-cartera}, después de presentar el diseño experimental y
la construcción de cartera del ganador.
